%!TeX root=../carlo_master.tex
\chapter{Computational resources}\label{app:compute}

The empirical campaigns of Sections~\ref{sec:binary_experiments} and~\ref{sec:multiclass_experiments} ran on rented cloud virtual machines; the synthetic studies of Section~\ref{sec:voronoi_comparison}, the boundary-rotation scenario of Section~\ref{sec:boundary_rotation_experiment}, and all table and figure generation ran locally on a personal laptop. Cloud machine-type selection, bootstrap, and run discipline are documented in \texttt{experiments/cloud/README.md}.

\paragraph{Cloud hardware.} Two virtual-machine types covered the OpenML benchmarks, running in parallel and merging their per-seed parquet shards on the laptop after each campaign. A Google Cloud CPU instance (\texttt{n2-highcpu-80}: 80 Intel Ice~Lake vCPUs, 80\,GB RAM, Ubuntu 24.04~LTS) ran the seven non-neural base estimators (GaussianNB, Logistic Regression, Random Forest, Extra Trees, XGBoost, LightGBM, CatBoost). A RunPod pod with one NVIDIA H200 (141\,GB HBM3e) on RunPod's CUDA~12.8 / PyTorch~2.8 PyTorch image ran the three deep tabular estimators (TabM, TabPFN, TabICLv2).

\paragraph{Per-study compute.} Each campaign uses ten replicate seeds and the dataset, model, strategy, and objective registries documented in its experimental-setup subsection. The merged parquets hold $136{,}680$ (dataset, model, objective, strategy, seed) rows for binary and $276{,}960$ for multiclass. Wall clock per leg, rounded: the CPU VM ran for ${\approx}\,10\,\text{min}$ on binary and ${\approx}\,2\,\text{h}$ on multiclass, and the H200 pod for ${\approx}\,2\,\text{h}$ on binary and ${\approx}\,25\,\text{min}$ on multiclass. The decision-stage \texttt{fit} calls themselves are a small fraction of total compute: the per-row \texttt{strategy\_fit\_seconds} field of the merged multiclass parquet sums to ${\approx}\,35\,\text{CPU-h}$ across all joblib workers, dominated by the five cross-validated tuning variants at ${\approx}\,5\text{--}6\,\text{CPU-h}$ each; on the binary side the harness caches base-model probabilities and replays them across objectives rather than refitting, keeping the recorded decision-stage cost negligible (${<}\,0.1\,\text{CPU-h}$).

\paragraph{Performance engineering.} The wall-clock figures above reflect a sequence of targeted optimizations to the harness and the \texttt{threshold\_opt} library; without them the multiclass campaign would have been several orders of magnitude more expensive. Three categories account for most of the gain. First, sharing cross-validation fold refits across the seven multiclass CV-variant strategies, which all use the same deterministic stratified $K$-fold split: refitting once per (dataset, model) and reusing through a pipeline-level cache dropped a representative smoke run from ${\approx}\,2.5\,\text{h}$ to under $30\,\text{min}$. Second, vectorizing the partition-optimizer inner loops and replacing finite-difference gradients in the Dirichlet calibrator with an analytical Jacobian, each yielding order-of-magnitude per-call speedups. Third, oversubscription control on the 80-vCPU CPU VM\@: base estimators are pinned to a single thread and the surrounding BLAS, OpenMP, and PyTorch pools are capped to one thread per joblib worker.

\paragraph{Total spend.} Total cloud spend over the project lifetime was approximately USD 140 on RunPod GPU and BRL 480 on Google Cloud CPU, inclusive of development-time smoke runs, GPU connection checks, and occasional resumed runs. The synthetic studies and the analysis pipeline added no cloud cost.

\paragraph{Reproducibility.} The merged parquets and per-row metadata (ten replicate seeds sampled deterministically from a master seed, the per-row \texttt{strategy\_fit\_seconds}, and the model-side test AUCs on both data partitions) are stored under \texttt{experiments/results/} and re-run end-to-end via \texttt{make experiments-binary} and \texttt{make experiments-multiclass}; the environment is pinned through a fixed \texttt{.python-version} and a committed \texttt{uv.lock}. A \texttt{pytest} suite over \texttt{src/} and \texttt{experiments/} (\texttt{make test}) covers the library and the harness, \texttt{make test-reproducibility} asserts that a fixed seed yields identical fits and predictions, and \texttt{make test-freshness} regenerates every committed table and figure from the stored results and fails on any diff against the versioned copy; the same suite runs in continuous integration on every push and pull request.
